\subsection{Adding a Multiple of One Row to Another}

Start with
\[
A=\begin{pmatrix}a&b\\c&d\end{pmatrix}
\]
and replace the second row by
\[
(c,d)+k(a,b)=(c+ka,d+kb).
\]

\begin{derivationbox}[The added terms cancel]
\begin{align*}
\det\begin{pmatrix}a&b\\c+ka&d+kb\end{pmatrix}
&=a(d+kb)-b(c+ka)\\
&=ad+kab-bc-kab\\
&=ad-bc\\
&=\det A.
\end{align*}
\end{derivationbox}

\begin{theorembox}[Row replacement]
Adding a multiple of one row to another row does not change the determinant.
\end{theorembox}

This can also be reconstructed from linearity: the added part creates a determinant with proportional rows, and that determinant is zero.

\begin{examplebox}[A determinant-preserving operation]
For
\[
A=\begin{pmatrix}2&1\\5&3\end{pmatrix},
\]
replace $R_2$ by $R_2-2R_1=(1,1)$. Both the original and transformed matrices have determinant $1$.
\end{examplebox}
